\section{The Return Passes Through the Neighbor}

If worship belongs to justice without equivalence, takes embodied common form in the Eucharist, and is safeguarded by prudent precept, one lethal counterfeit still haunts religious practice in every century: the attempt to offer God immaculate ritual at the altar while withholding justice and mercy from human beings outside the door. Scripture condemns that split as an abomination that renders outward worship a lie.

\subsection{The Prophetic and Gospel Refusal of Severed Worship}

Across the Old Testament, the prophets confront a worshiping people whose sacrifices are rubrically strict and morally violent. Through Isaiah, the Lord speaks with blunt revulsion against offerings brought by blood-stained hands:

\begin{quote}
To what purpose do you offer me the multitude of your victims, saith the Lord? I am full, I desire not holocausts of rams, and fat of fatlings, and blood of calves, and lambs, and buck goats. \dots\ And when you stretch forth your hands, I will turn away my eyes from you: and when you multiply prayer, I will not hear: for your hands are full of blood. Wash yourselves, be clean, take away the evil of your devices from my eyes, cease to do perversely, learn to do well: seek judgment, relieve the oppressed, judge for the fatherless, defend the widow. (Isa 1:11, 15--17)
\end{quote}

Amos delivers the same verdict against solemn assemblies divorced from righteousness: ``Take away from me the tumult of thy songs: and I will not hear the canticles of thy harp. But judgment shall be revealed as water, and justice as a mighty torrent'' (Amos 5:23--24). Micah poses the question of unpayable return---whether thousands of rams or even one's firstborn could appease the High God---and cuts straight to the substance: ``I will shew thee, O man, what is good, and what the Lord requireth of thee: Verily to do judgment, and to love mercy, and to walk solicitous with thy God'' (Mic 6:6--8). Hosea compresses the prophetic indictment into one sentence that Christ twice makes his own: ``For I desired mercy, and not sacrifice: and the knowledge of God more than holocausts'' (Hos 6:6; Matt 9:13; 12:7).

In the Sermon on the Mount, Christ carries this rule right up to the steps of the altar:

\begin{quote}
If therefore thou offer thy gift at the altar, and there thou remember that thy brother hath anything against thee; leave there thy offering before the altar, and go first to be reconciled to thy brother, and then coming thou shalt offer thy gift. (Matt 5:23--24)
\end{quote}

Worship is not abolished; the worshiper is told to return and offer his gift. Yet God refuses to accept a gift offered across an unhealed breach of fraternal justice. St.\ Paul enforces the same law at the Eucharistic table in Corinth, where wealthier Christians ate and drank sumptuously before the poor arrived and left them hungry and shamed: ``When you come therefore together into one place, it is not now to eat the Lord's supper\dots\ Or despise ye the church of God and put them to shame that have not?'' (1 Cor 11:20, 22). Failing to discern the Body of Christ at the altar and despising his poor members at the same gathering are two faces of one sacrilege.

And in the judgment scene of Matthew 25, Christ removes every remaining partition between devotion to him and care for the needy: ``For I was hungry, and you gave me to eat: I was thirsty, and you gave me to drink: I was a stranger, and you took me in\dots\ Amen I say to you, as long as you did it to one of these my least brethren, you did it to me'' (Matt 25:35, 40).

\subsection{Chrysostom on the Altar and the Naked Poor}

No Father of the Church drove this truth home with fiercer clarity than St.\ John Chrysostom. Preaching in Antioch on Matthew 14 (Homily 50), Chrysostom first leads his congregation to tremble with awe before the Real Presence in the Eucharist: Christ gives us not merely the hem of his garment to touch, but his whole sacrificed Body to receive, Christ's own hand stretching out through the priest's hand (\emph{Homilies on Matthew}, Hom.\ 50.3).\endnote{St.\ John Chrysostom, \emph{Homilies on the Gospel of St.\ Matthew}, Hom.\ 50.3--4 (Prevost--Riddle translation, NPNF I.10).}

Then, without lowering the reverence of the altar by an inch, Chrysostom turns that Eucharistic realism upon the conscience of the wealthy who adorn sanctuaries while ignoring the destitute:

\begin{quote}
Wouldest thou do honor to Christ's body? Neglect Him not when naked; do not while here thou honorest Him with silken garments, neglect Him perishing without of cold and nakedness. For He that said, ``This is my body,'' and by His word confirmed the fact, This same said, ``Ye saw me an hungered, and fed me not;'' and, ``Inasmuch as ye did it not to one of the least of these, ye did it not to me.'' For This indeed needs not coverings, but a pure soul; but that requires much attention.\endnote{Chrysostom, \emph{Homilies on Matthew}, Hom.\ 50.3 (NPNF I.10).}
\end{quote}

Chrysostom presses the contradiction with concrete irony. What would a starving man say if someone refused him bread while plating his table in silver, or left him shivering in rags while erecting golden columns ``to his honor''? He would count it mockery and extreme insult. ``Since God hath no need at all of golden vessels, but of golden souls,'' Chrysostom continues, ``do not therefore while adorning His house overlook thy brother in distress, for he is more properly a temple than the other'' (Hom.\ 50.3--4).

Yet Chrysostom refuses the false counter-extreme that would despise liturgical beauty in the name of charity. Twice in consecutive paragraphs he guards the Catholic balance explicitly:

\begin{quote}
And these things I say, not forbidding such offerings to be provided; but requiring you, together with them, and before them, to give alms. For He accepts indeed the former, but much more the latter. For in the one the offerer alone is profited, but in the other the receiver also. Here the act seems to be a ground even of ostentation; but there all is mercifulness, and love to man.\endnote{Chrysostom, \emph{Homilies on Matthew}, Hom.\ 50.3--4 (NPNF I.10).}
\end{quote}

Sacred vessels and beautiful churches remain fitting signs of reverence when offered from honest labor and a generous heart. What Christ rejects is the bribe: using sanctuary splendor as a cheap substitute for justice, or pretending to honor Christ on the altar while ignoring the very same Christ freezing in the street.

\subsection{Mercy as True Sacrifice: Irenaeus, Augustine, and Aquinas}

The Western Fathers and St.\ Thomas Aquinas explain the theological reason why care for the neighbor belongs to the very nature of sacrifice itself.

St.\ Irenaeus, contrasting the accepted offering of Abel with the rejected offering of Cain (\emph{Adversus Haereses} IV.18.3), insists that outward correctness cannot sanitize a heart divided by malice toward one's brother: ``Sacrifices, therefore, do not sanctify a man, for God stands in no need of sacrifice; but it is the conscience of the offerer that sanctifies the sacrifice when it is pure, and thus moves God to accept [the offering] as from a friend.''\endnote{Irenaeus, \emph{Adversus Haereses} IV.18.3 and IV.18.6 (ANF 1).} And in the closing paragraph of that chapter (IV.18.6), Irenaeus binds altar and almsgiving directly together through Proverbs 19:17 and Matthew 25: ``For even as God does not need our possessions, so do we need to offer something to God; as Solomon says: `He that hath pity upon the poor, lendeth unto the Lord.' For God, who stands in need of nothing, takes our good works to Himself for this purpose, that He may grant us a recompense of His own good things.''

St.\ Augustine unfolds the deepest synthesis of Hosea 6:6 in \emph{City of God} X.5--6.\endnote{Augustine, \emph{De civitate Dei} X.5--6 (Dods 1871; CSEL 40/1).} When God declares, ``I desire mercy rather than sacrifice,'' Augustine observes that he is not abolishing sacrifice, but preferring the true sacrifice to its outward sign:

\begin{quote}
And so, when it is written, ``I desire mercy rather than sacrifice,'' nothing else is meant than that one sacrifice is preferred to another; for that which in common speech is called sacrifice is only the symbol of the true sacrifice. Now mercy is the true sacrifice, and therefore it is said, as I have just quoted, ``with such sacrifices God is well pleased.''\endnote{Augustine, \emph{De civitate Dei} X.5 (Dods 1871; quoting Heb 13:16; CSEL 40/1: \emph{Porro autem misericordia uerum sacrificium est}).}
\end{quote}

In chapter 6, Augustine adds an essential precision: ``And therefore even the mercy we show to men, if it is not shown for God's sake, is not a sacrifice.'' Philanthropy cut off from God may relieve a temporal hardship, yet it does not reach the level of sacrifice unless it is referred to God as the supreme good in whom alone human misery is finally healed. Conversely, when works of mercy---both toward the neighbor in distress and toward our own soul in repentance and temperance---are done for God's sake, they become part of the universal sacrifice offered by the whole redeemed City through Christ the High Priest.

St.\ Thomas Aquinas brings these patristic insights into clear scholastic order. Explaining St.\ James's definition of ``religion clean and undefiled'' (Jas 1:27), he distinguishes between the virtue that elicits an act and the virtue that directs it: visiting orphans and widows in their affliction is elicited by the virtue of mercy, yet \emph{commanded} by the virtue of religion when directed to God's honor. Turning to sacrifice, he adds that external goods are offered to God in two ways: ``directly when we offer our possession to God immediately, and indirectly when we share them with our neighbor for God's sake.''\endnote{Aquinas, \emph{ST} II-II, q.\ 81, a.\ 1, ad 1; II-II, q.\ 85, a.\ 3, ad 2.}

Here scholastic theology rejoins St.\ Justin Martyr's Sunday assembly (\emph{First Apology} 67). The bread and wine brought to the altar and the alms placed with the president for orphans, widows, the sick, prisoners, and strangers belong to one movement of gratitude: because God needs nothing for himself, the Creator who receives nothing into his own fullness redirects our return outward into his needy members.
